% part: intuitionistic-logic
% chapter: soundness-completeness
% section: canonical-model

\documentclass[../../../include/open-logic-section]{subfiles}

\begin{document}

\olfileid[pt]{int}{sc}{mod}

\olsection{O modelo canônico}

Os mundos de nosso modelo serão sequências finitas~$\sigma$ de números
naturais, O.s.\@, $\sigma \in \Nat^*$. Observe que $\Nat^*$ é definido
indutivamente por:
\begin{enumerate}
\item $\emptyseq \in \Nat^*$.
\item Se $\sigma \in \Nat^*$ e $n \in \Nat$, então $\sigma.n \in
  \Nat^*$ (em que $\sigma.n$ é $\sigma \concat \tuple{n}$ e
  $\sigma \concat \sigma'$ é a concatenação de $\sigma$ e $\sigma'$).
\item Nada mais pertence a $\Nat^*$.
\end{enumerate}
Assim, podemos usar $\Nat^*$ para formular definições indutivas.

Seja $\tuple{!B_1, !C_1}$, $\tuple{!B_2, !C_2}$, \dots, uma
enumeração de todos os pares de !!{formula}s. Dado um conjunto de
!!{formula}s~$\Delta$, defina $\Delta(\sigma)$ por indução da seguinte
maneira:
\begin{enumerate}
\item $\Delta(\emptyseq) = \Delta$
\item $\Delta(\sigma.n) = {}$
  \[
  \begin{cases}
    (\Delta(\sigma) \cup \{!B_n\})^* &
    \text{se $\Delta(\sigma) \cup \{!B_n\} \Proves/ !C_n$} \\
    \Delta(\sigma) & \text{caso contrário}
  \end{cases}
  \]
\end{enumerate}
Por $(\Delta(\sigma) \cup \{!B_n\})^*$ entendemos aqui o conjunto
primo de !!{formula}s cuja existência é garantida por
\olref[lin]{lem:lindenbaum}, aplicada ao conjunto
$\Delta(\sigma) \cup \{!B_n\}$ e à !!{formula}~$!C_n$. Observe que, por
essa definição, se $\Delta(\sigma) \cup \{!B_n\} \Proves/ !C_n$,
então $\Delta(\sigma.n) \Proves !B_n$ e
$\Delta(\sigma.n) \Proves/ !C_n$. Observe também que
$\Delta(\sigma) \subseteq \Delta(\sigma.n)$ para todo~$n$. Se
$\Delta$ é primo, então $\Delta(\sigma)$ é primo para todo~$\sigma$.

\begin{defn}\ollabel{defn:canonical-model}
  Suponha que $\Delta$ seja primo. Então o \emph{modelo canônico}
  $\mModel{M(\Delta)}$ para $\Delta$ é definido por:
  \begin{enumerate}
  \item $W = \Nat^*$, o conjunto das sequências finitas de números naturais.
  \item $R$ é a ordem parcial segundo a qual $R\sigma\sigma'$ sse
    $\sigma$ é um segmento inicial de~$\sigma'$ (O.s.\@,
    $\sigma' = \sigma \concat \sigma''$ para alguma sequência~$\sigma''$).
  \item $V(p) = \Setabs{\sigma}{p \in \Delta(\sigma)}$.
  \end{enumerate}
\end{defn}

É fácil verificar que $R$ é de fato uma ordem parcial. Além disso, a
condição de monotonicidade sobre~$V$ é satisfeita. Como
$\Delta(\sigma) \subseteq \Delta(\sigma.n)$, obtemos
$\Delta(\sigma) \subseteq \Delta(\sigma')$ sempre que
$R\sigma\sigma'$, por indução sobre~$\sigma$.

\end{document}
